%!TEX root = ../lecture_2.tex

%\titlegraphic {
%    \begin{tikzpicture}[overlay,remember picture]
%        \node[right=0.2cm] at (current page.150){\includegraphics[width=3cm]{Beamer-Logo}
%            \node[left=0.2cm] at (current page.30){\includegraphics[width=3cm]{Beamer-Logo}
%            };
%    \end{tikzpicture}
%}


\newcommand*{\LogoLeft}{
    \logo{
        \begin{tikzpicture}[overlay,remember picture]
            \node[right=0.2cm] at (current page.155){\includegraphics[width=1.5cm]{./images/logo_white.png}};
            \node[left=0.2cm] at (current page.26){\includegraphics[width=2.5cm]{./images/department_logo_white.png}};
        \end{tikzpicture}
    }
}

\newcommand*{\LogoRight}{
    \logo{
        \begin{tikzpicture}[overlay,remember picture]
            \node[left=0.2cm] at (current page.26){\includegraphics[width=1.0cm]{./images/logo_black.png}};
            \node[left=0.2cm] at (current page.335){\includegraphics[width=2.7cm]{./images/department_logo_black.png}};
        \end{tikzpicture}

    }
}

% Возможность использовать автонумерацию
\newcounter{index}
\newcommand*{\No}{\noindent\refstepcounter{index}\textnumero~\theindex}
\newcommand*{\Index}{\refstepcounter{index}\theindex}

%%% Русская традиция начертания математических знаков
\renewcommand{\le}{\ensuremath{\leqslant}}
\renewcommand{\leq}{\ensuremath{\leqslant}}
\renewcommand{\ge}{\ensuremath{\geqslant}}
\renewcommand{\geq}{\ensuremath{\geqslant}}
\renewcommand{\emptyset}{\varnothing}

%%% Русская традиция начертания математических функций
\renewcommand{\tan}{\operatorname{tg}}
\renewcommand{\cot}{\operatorname{ctg}}
\renewcommand{\csc}{\operatorname{cosec}}
%%% Русская традиция начертания греческих букв (пакет upgreek)
\renewcommand{\epsilon}{\ensuremath{\upvarepsilon}}   %  русская традиция записи
\renewcommand{\phi}{\ensuremath{\upvarphi}}
\renewcommand{\kappa}{\ensuremath{\varkappa}}
\renewcommand{\alpha}{\ensuremath{\upalpha}}
\renewcommand{\beta}{\ensuremath{\upbeta}}
\renewcommand{\gamma}{\ensuremath{\upgamma}}
\renewcommand{\delta}{\ensuremath{\updelta}}
\renewcommand{\varepsilon}{\ensuremath{\upvarepsilon}}
\renewcommand{\zeta}{\ensuremath{\upzeta}}
\renewcommand{\eta}{\ensuremath{\upeta}}
\renewcommand{\theta}{\ensuremath{\uptheta}}
\renewcommand{\vartheta}{\ensuremath{\upvartheta}}
\renewcommand{\iota}{\ensuremath{\upiota}}
\renewcommand{\kappa}{\ensuremath{\upkappa}}
\renewcommand{\lambda}{\ensuremath{\uplambda}}
\renewcommand{\mu}{\ensuremath{\upmu}}
\renewcommand{\nu}{\ensuremath{\upnu}}
\renewcommand{\xi}{\ensuremath{\upxi}}
\renewcommand{\pi}{\ensuremath{\uppi}}
\renewcommand{\varpi}{\ensuremath{\upvarpi}}
\renewcommand{\rho}{\ensuremath{\uprho}}
\renewcommand{\varrho}{\ensuremath{\upvarrho}}
\renewcommand{\sigma}{\ensuremath{\upsigma}}
\renewcommand{\varsigma}{\ensuremath{\upvarsigma}}
\renewcommand{\tau}{\ensuremath{\uptau}}
\renewcommand{\upsilon}{\ensuremath{\upupsilon}}
\renewcommand{\varphi}{\ensuremath{\upvarphi}}
\renewcommand{\chi}{\ensuremath{\upchi}}
\renewcommand{\psi}{\ensuremath{\uppsi}}
\renewcommand{\omega}{\ensuremath{\upomega}}
